% Qualitative illustration of Pyatakov lecture 12, slide at 42 min.
% Not a numerical solution or experimental trajectory.
\begin{tikzpicture}[>=Stealth,font=\small]
\begin{scope}[x=.95cm,y=1cm]
 \draw[->] (0,0)--(3.8,0) node[right] {$V$};
 \draw[->] (0,0)--(0,2.7) node[above] {$I_d$};
 \draw[thick,MidnightBlue] (0,0)..controls(.35,.4) and(.45,2.3)..(.9,2.3);
 \draw[thick,BrickRed] (.9,2.3)..controls(1.5,2.3) and(1.75,.65)..(2.4,.65);
 \draw[thick,MidnightBlue] (2.4,.65)..controls(3,.65) and(3.1,1.4)..(3.4,2.3);
 \node[align=center,font=\footnotesize] at (1.8,-.65) {На падающем участке\\$\dd I_d/\dd V<0$};
\end{scope}
\begin{scope}[shift={(7.2cm,1.3cm)},x={(.75cm,-.26cm)},y={(.68cm,.28cm)},z={(0cm,.9cm)}]
 \foreach \h in {-1.15,1.15}{
  \draw[MidnightBlue!35,fill=MidnightBlue!4] (-1.8,-1.4,\h)--(1.8,-1.4,\h)--(1.8,1.4,\h)--(-1.8,1.4,\h)--cycle;
  \draw[MidnightBlue,thick,domain=0:900,samples=170,variable=\t]
   plot ({(.15+.00105*\t)*cos(\t)},{(.15+.00105*\t)*sin(\t)},\h);
 }
 \draw[BrickRed,thick,->] (1.3,0,1.15) ..controls (2,-.2,.2) and (-1.8,.1,-.3)..(-.2,.2,-1.15);
 \draw[BrickRed,thick,->] (-1.25,0,-1.15) ..controls (-2,.2,-.3) and (1.8,.2,.4)..(.1,-.1,1.15);
 \draw[->] (0,0,-1.8)--(2.5,0,-1.8) node[right] {$q$};
 \draw[->] (0,0,-1.8)--(0,2.3,-1.8) node[right] {$I_L$};
 \draw[->] (0,0,-1.8)--(0,0,2) node[above] {$V$};
 \node at (1.6,1.2,1.15) {A};
 \node at (1.6,1.2,-1.15) {B};
\end{scope}
\end{tikzpicture}
